\chapter{Dynamics of Live Black-Hole Phase Boundaries}
\label{app:live_black_hole_boundaries}

\section*{Surface Motion, Spin, Binary Merger, and Ringdown in the FCC/SC Transition}

As established in TFT, a black hole is not a point singularity and not an empty geometric hole. It is a physical phase transition of the Spacetime Gravitational Floor. Under gravitational saturation, the active Face-Centered Cubic Soft Floor loses its interstitial freedom and locks into the Simple Cubic Hard Floor.

The event horizon is therefore interpreted as the physical FCC/SC phase boundary of the vacuum:
\begin{equation}
    \mathrm{FCC\ Soft\ Floor} \rightarrow \mathrm{SC\ phase\text{-}locked\ Hard\ Floor}.
\end{equation}

This appendix develops the next level of the theory: the dynamics of that boundary.

In TFT, a black hole is a compact saturated phase-domain with center-of-mass motion, orbital motion, axial angular momentum, surface growth, merger behavior, and relaxation modes. These dynamics are not assigned to ordinary material flow through the saturated SC interior. The SC domain is phase-locked, void-less, and mechanically saturated; it does not support ordinary internal circulation, interstitial transport, or fluid-like material rotation. The dynamical degrees of freedom instead reside on the live FCC/SC phase boundary and in the strained Soft Floor surrounding it.

A TFT black hole is therefore represented mechanically as:
\begin{equation}
    \text{Black Hole} = \text{SC Saturated Domain} + \text{FCC/SC Phase Boundary} + \text{External Strained Soft Floor}.
\end{equation}

This definition is local and mechanical. It does not depend on optical appearance, redshift interpretation, coordinate convention, or what a distant observer sees. It depends on the saturation state of the vacuum itself. Wherever the local gravitational saturation condition is reached,
\begin{equation}
    \alpha_{\mathrm{G}} = 1,
\end{equation}
the active FCC vacuum can no longer remain dynamically available and becomes part of the phase-locked SC saturated domain.

The radial black-hole chapters treated this process in the spherical average, where the saturation radius evolves according to:
\begin{equation}
    \dot{R}_{\mathrm{QG}} = \frac{G}{c_{\mathrm{E}}^2}\dot{M}.
\end{equation}

That relation describes average radial growth. The present appendix extends the same phase-boundary logic to nonspherical, rotating, and binary systems. The central task is to replace a purely radial black-hole picture with a live surface theory containing distinct but coupled channels: center-of-mass motion, phase-front growth, tangential redistribution, torsional spin, saturated bridge formation, stress shedding, and ringdown.

\section{Velocity Channels of a Live Black-Hole Boundary}

A dynamic black hole contains more than one form of motion. In the TFT description, these motions are assigned to distinct mechanical channels of the FCC/SC phase boundary:

\begin{itemize}
    \item \textbf{Center-of-Mass Motion:} The black hole behaves externally as a compact saturated phase-boundary object moving through the surrounding Soft Floor. This is the motion relevant to binary black-hole orbits, inspiral trajectories, merger recoil, and motion through an external gravitational field, describing the displacement of the entire saturated black-hole domain.
    \item \textbf{Phase-Front Motion:} This is the local advancement of the FCC/SC transition surface as the black hole gains mass-energy and converts surrounding active Soft Floor into the phase-locked SC substrate state: 
    $\text{FCC Soft Floor} \rightarrow \text{SC phase-locked substrate}$. 
    In the spherical growth limit, this channel is represented by the radial growth law $\dot{R}_{\mathrm{QG}} = \frac{G}{c_{\mathrm{E}}^2}\dot{M}$. In a disturbed or merging black hole, this conversion can occur unevenly, producing local bulging, distortion, and subsequent relaxation.
    \item \textbf{Tangential Redistribution:} Since the effective inertial mass of the black hole is phase-segregated to the boundary layer, surface mass density and surface stress may redistribute tangentially across the phase boundary. This occurs during matter inflow, tidal distortion, asymmetric growth, binary interaction, or post-merger relaxation, acting as a boundary-layer process rather than bulk transport.
    \item \textbf{Torsional Boundary Motion:} This is the channel associated with axial angular momentum. The observational angular momentum is carried mechanically by torsional loading of the FCC/SC boundary and shear strain in the surrounding Soft Floor.
\end{itemize}

Thus, the live boundary contains four distinct velocity channels:
\begin{equation}
    \text{center-of-mass motion} \neq \text{phase-front growth} \neq \text{tangential redistribution} \neq \text{torsional boundary motion}.
\end{equation}

These channels are distinct but not isolated. In a binary merger, rapid orbital motion can drive boundary deformation, bridge saturation, gravitational-wave stress shedding, and angular-momentum transfer. However, orbital motion is not itself the same mechanical process as phase-front growth or torsional boundary motion.

A black hole may therefore translate through the Soft Floor, grow by FCC-to-SC conversion, redistribute boundary stress, and carry axial angular momentum through different but coupled mechanical channels. This decomposition provides the basis for treating inspiral, spin, merger, gravitational-wave emission, and ringdown within the live FCC/SC phase-boundary framework.

\section{Axial Rotation as Boundary Torsion}

A rotating black hole carries axial angular momentum. In the TFT boundary description, this angular momentum is not assigned to ordinary material rotation of the SC saturated interior. State II is phase-locked; the saturated SC domain supplies the internal boundary condition, while the live FCC/SC phase boundary and the surrounding Soft Floor carry the rotational degree of freedom.

The mechanical interpretation is therefore

\begin{equation}
    \text{black-hole spin}
    =
    \text{boundary torsion}
    +
    \text{external Soft Floor shear}.
\end{equation}

Thus, a rotating TFT black hole is modeled as a saturated phase-domain whose boundary carries angular momentum and transmits that angular momentum into the surrounding vacuum as torsional strain.

For a first-order surface-shell approximation, treat the black-hole boundary as a thin mass-bearing layer. Let \(dM_{\partial}\) denote the boundary mass element and let

\begin{equation}
    r_\perp
    =
    R_{\mathrm{QG}}\sin\theta
\end{equation}

be the cylindrical distance from the spin axis. The surface moment of inertia is

\begin{equation}
    I_{\mathrm{surf}}
    =
    \int_{\partial_{\mathrm{BH}}}
    r_\perp^2\,dM_{\partial},
\end{equation}

where \(\partial_{\mathrm{BH}}\) denotes the live black-hole phase boundary.

The corresponding axial angular momentum is, in general,

\begin{equation}
    \mathcal{J}_{\mathrm{BH}}
    =
    \int_{\partial_{\mathrm{BH}}}
    r_\perp\,v_{\phi}\,dM_{\partial},
\end{equation}

where \(v_{\phi}\) is the azimuthal boundary speed. For a coherent axisymmetric torsional mode satisfying

\begin{equation}
    v_{\phi}
    =
    \frac{v_{\mathrm{tor}}}{R_{\mathrm{QG}}}r_\perp,
\end{equation}

this reduces to

\begin{equation}
    \mathcal{J}_{\mathrm{BH}}
    =
    I_{\mathrm{surf}}
    \frac{v_{\mathrm{tor}}}{R_{\mathrm{QG}}},
\end{equation}

where \(v_{\mathrm{tor}}\) is the effective torsional boundary speed, evaluated at the boundary scale.

The dimensionless spin parameter is

\begin{equation}
    a_\ast
    =
    \frac{c_{\mathrm{E}}\mathcal{J}_{\mathrm{BH}}}{GM^2}.
\end{equation}

Substituting the surface-shell expression gives

\begin{equation}
    a_\ast
    =
    \frac{c_{\mathrm{E}}}{GM^2}
    I_{\mathrm{surf}}
    \frac{v_{\mathrm{tor}}}{R_{\mathrm{QG}}}.
\end{equation}

Using the TFT saturation radius,

\begin{equation}
    R_{\mathrm{QG}}
    =
    \frac{GM}{c_{\mathrm{E}}^2},
\end{equation}

the spin parameter becomes

\begin{equation}
    a_\ast
    =
    \left(
    \frac{I_{\mathrm{surf}}}{M R_{\mathrm{QG}}^2}
    \right)
    \frac{v_{\mathrm{tor}}}{c_{\mathrm{E}}}.
\end{equation}

The dimensionless factor

\begin{equation}
    \frac{I_{\mathrm{surf}}}{M R_{\mathrm{QG}}^2}
\end{equation}

is the normalized surface moment of inertia of the live boundary. It is not an independent constant; it is determined by the boundary mass-current distribution.

For a uniformly loaded spherical boundary shell,

\begin{equation}
    I_{\mathrm{surf}}
    =
    \frac{2}{3}
    M R_{\mathrm{QG}}^2,
\end{equation}

and therefore

\begin{equation}
    a_\ast
    =
    \frac{2}{3}
    \frac{v_{\mathrm{tor}}}{c_{\mathrm{E}}}.
\end{equation}

For an equatorially concentrated boundary current,

\begin{equation}
    I_{\mathrm{surf}}
    \rightarrow
    M R_{\mathrm{QG}}^2,
\end{equation}

and therefore

\begin{equation}
    a_\ast
    \rightarrow
    \frac{v_{\mathrm{tor}}}{c_{\mathrm{E}}}.
\end{equation}

For any boundary-confined mass distribution,

\begin{equation}
    0
    \leq
    r_\perp
    \leq
    R_{\mathrm{QG}},
\end{equation}

so

\begin{equation}
    I_{\mathrm{surf}}
    \leq
    M R_{\mathrm{QG}}^2.
\end{equation}

The torsional boundary speed is further constrained by the finite kinematic capacity of the Soft Floor,

\begin{equation}
    v_{\mathrm{tor}}
    \leq
    c_{\mathrm{E}}.
\end{equation}

Therefore, within the first-order boundary-confined surface-shell approximation,

\begin{equation}
    a_\ast
    \leq
    1.
\end{equation}

In TFT, this spin bound is interpreted as a boundary-kinematic saturation condition. The phase boundary cannot carry torsional motion faster than the elastic response capacity of the Soft Floor, and the effective angular-momentum carrier is confined to the black-hole boundary scale.

This result identifies the mechanical carrier of black-hole angular momentum in the live-boundary model. It is not a full Kerr replacement. A Kerr-equivalent TFT solution requires the additional derivation of the external torsional Soft-Floor field, the frame-dragging analogue, the spin-dependent boundary geometry, the horizon-scale mapping, and the corresponding strong-field orbital dynamics.

\section{Binary Black-Hole Orbital Motion}

A black hole can participate in orbital motion because, externally, it behaves as a compact saturated phase-boundary object embedded in the surrounding Soft Floor. Its center-of-mass motion responds to the strain field of another compact object in the same general way that any localized mass-defect responds to an external gravitational field. This orbital motion is distinct from phase-front growth, tangential boundary redistribution, and torsional boundary spin, although these channels can become strongly coupled during merger.

For an ideal circular binary with component masses $M_1$ and $M_2$, define the total mass:
\begin{equation}
    M_{\mathrm{tot}} = M_1 + M_2.
\end{equation}

Let $d$ be the center-to-center separation between the two saturated domains. The relative orbital velocity scale is:
\begin{equation}
    v_{\mathrm{rel}}^2 = \frac{G M_{\mathrm{tot}}}{d}.
\end{equation}

This is the relative speed between the two black holes in the two-body frame. The individual orbital speeds about the center of mass are:
\begin{equation}
    v_1 = \frac{M_2}{M_{\mathrm{tot}}}v_{\mathrm{rel}},
\end{equation}
and
\begin{equation}
    v_2 = \frac{M_1}{M_{\mathrm{tot}}}v_{\mathrm{rel}}.
\end{equation}

Each TFT black-hole saturation radius is:
\begin{equation}
    R_i = \frac{G M_i}{c_{\mathrm{E}}^2}.
\end{equation}

Therefore,
\begin{equation}
    R_1 + R_2 = \frac{G(M_1+M_2)}{c_{\mathrm{E}}^2} = \frac{G M_{\mathrm{tot}}}{c_{\mathrm{E}}^2}.
\end{equation}

Substituting this directly into the relative-velocity expression gives:
\begin{equation}
    v_{\mathrm{rel}}^2 = c_{\mathrm{E}}^2 \left( \frac{R_1+R_2}{d} \right).
\end{equation}

Therefore,
\begin{equation}
    v_{\mathrm{rel}} = c_{\mathrm{E}} \sqrt{\frac{R_1+R_2}{d}}.
\end{equation}

Here, the controlling quantity is the physical ratio $\frac{R_1+R_2}{d}$.

When the two saturation boundaries are widely separated:
\begin{equation}
    d > R_1 + R_2 \quad \Longrightarrow \quad v_{\mathrm{rel}} < c_{\mathrm{E}}.
\end{equation}

In the ideal circular boundary-contact limit:
\begin{equation}
    d = R_1 + R_2 \quad \Longrightarrow \quad v_{\mathrm{rel}} = c_{\mathrm{E}}.
\end{equation}

For equal masses, $M_1 = M_2$, the individual speeds about the center of mass are $v_1 = v_2 = \frac{1}{2}v_{\mathrm{rel}}$. At the ideal contact limit:
\begin{equation}
    v_1 = v_2 = \frac{c_{\mathrm{E}}}{2}.
\end{equation}

Therefore, the $c_{\mathrm{E}}$ condition is not a universal binary velocity law. It is a limiting contact-scale result for the relative motion of an ideal circular binary. The more general relation is controlled by the separation ratio:
\begin{equation}
    v_{\mathrm{rel}} = c_{\mathrm{E}} \sqrt{\frac{R_1+R_2}{d}}.
\end{equation}

In the TFT phase-boundary picture, the physical merger need not wait for exact geometric contact of the two saturation radii. Merger may begin earlier if the strained FCC region between the two black holes reaches saturation and forms a common SC bridge. Thus, the contact-limit velocity should be treated as a useful limiting scale, not as a required velocity for all binary black-hole mergers.

This relation may later be formalized as a TFT contact-limit statement:
\begin{equation}
    d \rightarrow R_1 + R_2 \quad \Longrightarrow \quad v_{\mathrm{rel}} \rightarrow c_{\mathrm{E}},
\end{equation}
but only under the assumptions of circular motion, center-to-center separation, and ideal boundary contact. It should not be overstated as a general rule for arbitrary inspiral, eccentric motion, unequal-mass merger, or nonlinear phase-boundary coalescence.

\section{Merger as Saturated Bridge Formation}

The physical merger of two TFT black holes is not the collision of two rigid bodies, and it is not the geometric limit $d=0$. It is the formation of a common saturated phase boundary between two previously separate FCC/SC transition domains.

As two black holes inspiral, their external FCC strain fields increasingly overlap. The bridge region between the two saturated domains becomes compressed, sheared, and dynamically loaded by orbital motion, tidal distortion, boundary deformation, and spin-induced torsion. When the local Soft Floor in this bridge region can no longer remain in the active FCC state, it undergoes saturation and locks into the SC phase.

The merger therefore begins when the intermediate FCC region reaches the same phase condition that defines a black-hole boundary:
\begin{equation}
    \text{FCC Soft Floor} \rightarrow \text{SC phase-locked substrate}.
\end{equation}

At that point, the two initially separate boundaries are no longer independent. A common FCC/SC transition surface forms, and the system becomes a single saturated phase-boundary domain.

The merger sequence may be written schematically as:
\begin{equation}
\begin{split}
    \text{two separate SC domains} &\rightarrow \text{overlapping FCC strain fields} \rightarrow \text{saturated bridge} \\
    &\rightarrow \text{common phase boundary} \rightarrow \text{single larger SC domain} \\
    &\rightarrow \text{surface relaxation/ringdown}.
\end{split}
\end{equation}

A compact provisional expression for the bridge condition is:
\begin{equation}
    \alpha_{\mathrm{bridge}} = \alpha_{\mathrm{G},1} + \alpha_{\mathrm{G},2} + \alpha_{\mathrm{dyn}} \geq 1.
\end{equation}

Here $\alpha_{\mathrm{G},1}$ and $\alpha_{\mathrm{G},2}$ are the baseline gravitational strain contributions from the two black holes at the bridge location. The additional term $\alpha_{\mathrm{dyn}}$ represents dynamic loading of the bridge region due to orbital compression, shear, torsional boundary motion, tidal deformation, and nonlinear phase-boundary distortion.

More locally, let $\mathcal{B}_{\mathrm{BH12}}$ denote the bridge region between the two saturated black-hole domains. A common boundary begins to form when there exists a connected saturated bridge subset $\mathcal{B}_{\mathrm{sat}}(t) \subset \mathcal{B}_{\mathrm{BH12}}$ such that
\begin{equation}
    \alpha_{\mathrm{bridge}}(\mathbf{x},t) \geq 1 \qquad \forall\,\mathbf{x} \in \mathcal{B}_{\mathrm{sat}}(t),
\end{equation}
and $\mathcal{B}_{\mathrm{sat}}(t)$ connects the two initially separate FCC/SC boundaries.

This means that merger is triggered by the formation of a connected saturated bridge in the intervening Soft Floor, not by zero separation of the two centers. The common boundary forms when the bridge region crosses the FCC-to-SC phase threshold and becomes part of a shared phase-locked substrate domain.

The term $\alpha_{\mathrm{dyn}}$ should be read as a placeholder for future strong-field surface dynamics. It is not yet a closed-form field equation. Its role in this appendix is to mark the fact that binary merger cannot be determined from the two quasi-steady strain contributions alone once the system enters the nonlinear, rapidly evolving phase-boundary regime.

Thus, the TFT merger criterion is a local mechanical saturation condition:
\begin{equation}
    \boxed{ \text{black-hole merger} = \text{formation of a saturated SC bridge and common FCC/SC boundary} }.
\end{equation}

This definition is distinct from visual contact, rigid-body collision, or the formal limit $d \rightarrow 0$. In TFT, the physical merger begins when the Soft Floor between the two black holes loses its ability to remain an active FCC bridge and becomes part of the shared phase-locked substrate domain.

\section{Surface Growth and Mass Redistribution}

The radial growth law derived in \hyperref[app:phase_transition]{Appendix~\ref*{app:phase_transition}} describes the average spherical expansion of the saturation boundary:
\begin{equation}
    \dot{R}_{\mathrm{QG}} = \frac{G}{c_{\mathrm{E}}^2}\dot{M}.
\end{equation}

This relation is sufficient for a spherically averaged black hole, but a live boundary requires a local surface description. During accretion, spin loading, tidal distortion, or merger, mass-energy and strain need not be deposited uniformly over the boundary. Some regions of the FCC/SC transition surface may advance faster than others, producing local bulging, flattening, distortion, and subsequent relaxation.

A minimally deformed boundary may be written schematically as:
\begin{equation}
    R(\theta,\varphi,t) = R_{\mathrm{QG}}(t) + \delta R(\theta,\varphi,t),
\end{equation}
where $R_{\mathrm{QG}}(t)$ is the average saturation radius and $\delta R(\theta,\varphi,t)$ is the local boundary deformation.

The normal boundary speed is:
\begin{equation}
    v_n = \frac{\partial R}{\partial t}.
\end{equation}
Here $v_n$ measures the local normal motion of the FCC/SC phase front. Positive $v_n$ corresponds to local phase-front advancement. Local relaxation may reduce $\delta R$ relative to the average boundary shape, but this should not be interpreted as global black-hole evaporation or reversal of the SC phase unless a separate TFT mechanism is introduced.

To describe nonuniform boundary loading, define the effective surface mass density as:
\begin{equation}
    \Sigma_{\partial} = \frac{\mathrm{d}M_{\partial}}{\mathrm{d}A},
\end{equation}
where $M_{\partial}$ denotes the mass-energy dynamically encoded on the boundary layer to track the boundary-segregated inertial content of the black hole during nonspherical growth, redistribution, and relaxation.

Surface mass changes through four mechanical channels:
\begin{itemize}
    \item \textbf{Incoming Matter-Energy:} $\dot{\Sigma}_{\partial,\mathrm{in}}$, the surface mass-energy deposited by infalling matter or radiation.
    \item \textbf{Vacuum Phase Conversion:} $\dot{\Sigma}_{\partial,\mathrm{vac}}$, the mass-equivalent contribution associated with FCC-to-SC vacuum phase conversion at the boundary.
    \item \textbf{Tangential Redistribution:} $\nabla_A\left(\Sigma_{\partial} v_{\parallel}^{A}\right)$, describing tangential redistribution of boundary mass and stress along the live phase surface.
    \item \textbf{Stress Shedding:} $\dot{\Sigma}_{\partial,\mathrm{shed}}$, the mass-equivalent energy shed from the boundary as gravitational-wave emission, torsional strain radiation, or other stress-relaxation channels.
\end{itemize}

For a near-spherical boundary, a schematic conservation form is:
\begin{equation}
    \frac{\partial \Sigma_{\partial}}{\partial t} + \nabla_A\left(\Sigma_{\partial} v_{\parallel}^{A}\right) + \frac{2\Sigma_{\partial}v_n}{R} = \dot{\Sigma}_{\partial,\mathrm{in}} + \dot{\Sigma}_{\partial,\mathrm{vac}} - \dot{\Sigma}_{\partial,\mathrm{shed}}.
\end{equation}

The expansion term $\frac{2\Sigma_{\partial}v_n}{R}$ represents the dilution of surface density caused by local area growth in the near-spherical approximation. If the boundary area expands while the same mass is spread over a larger surface, the surface density decreases unless compensated by additional inflow or phase conversion.

This equation should be read as a boundary bookkeeping relation, not as a completed fundamental field equation. Its role is to state how the live black-hole surface can grow, redistribute, deform, and shed stress while remaining consistent with the averaged radial law $\dot{R}_{\mathrm{QG}} = \frac{G}{c_{\mathrm{E}}^2}\dot{M}$.

In the spherical limit,
\begin{equation}
    \delta R \rightarrow 0, \qquad v_{\parallel}^{A} \rightarrow 0, \qquad \Sigma_{\partial} \rightarrow \frac{M}{4\pi R_{\mathrm{QG}}^2},
\end{equation}
the local description collapses back to the averaged growth model.

\section{Stress Shedding and Ringdown}

When two black holes merge, the newly formed common FCC/SC boundary is not immediately relaxed. It is distorted, dynamically loaded, and torsionally strained. The live boundary must therefore shed excess energy, redistribute angular momentum, and settle toward a lower-strain saturated configuration.

The final mass of the remnant may be written as:
\begin{equation}
    M_f = M_1 + M_2 - \frac{E_{\mathrm{shed}}}{c_{\mathrm{E}}^2},
\end{equation}
where $E_{\mathrm{shed}}$ is the energy carried away from the merger system through gravitational-wave strain, torsional Soft Floor radiation, and other boundary-relaxation channels.

Thus, the final axial angular momentum may be written as:
\begin{equation}
    \mathcal{J}_f = \mathcal{J}_1^{(\mathrm{spin})} + \mathcal{J}_2^{(\mathrm{spin})} + \mathcal{J}_{\mathrm{orb}}^{(\mathrm{bin})} - \mathcal{J}_{\mathrm{shed}}.
\end{equation}

Here $\mathcal{J}_1^{(\mathrm{spin})}$ and $\mathcal{J}_2^{(\mathrm{spin})}$ are the initial spin angular momenta of the two black holes, $\mathcal{J}_{\mathrm{orb}}^{(\mathrm{bin})}$ is the binary orbital angular momentum transferred into the merger, and $\mathcal{J}_{\mathrm{shed}}$ is the angular momentum radiated away or redistributed through outgoing Soft Floor strain.

In TFT, gravitational waves from black-hole merger are interpreted as stress-shedding waves of the live phase boundary. They are produced when the common FCC/SC transition surface rapidly reconfigures and exports excess shear, compression, and torsional loading into the surrounding Soft Floor.

The ringdown stage is the final relaxation of the distorted phase boundary. Let the post-merger boundary be written as:
\begin{equation}
    R(\theta,\varphi,t) = R_{\mathrm{QG},f} + \delta R(\theta,\varphi,t),
\end{equation}
where $R_{\mathrm{QG},f}$ is the average saturation radius of the final remnant and $\delta R$ is the residual boundary deformation.

A minimal linearized relaxation model is:
\begin{equation}
    \frac{\partial^2 \delta R}{\partial t^2} + \Gamma_{\mathrm{QG}}\frac{\partial \delta R}{\partial t} + \omega_{\mathrm{QG}}^2\delta R = 0.
\end{equation}

Here $\Gamma_{\mathrm{QG}}$ is the effective damping rate of the post-merger boundary deformation, and $\omega_{\mathrm{QG}}$ is the natural relaxation frequency of the live FCC/SC boundary.
The physical interpretation is:
\begin{equation}
    \text{ringdown} = \text{damped relaxation of the live FCC/SC phase boundary}.
\end{equation}

This keeps the TFT description close to observation: the post-merger signal remains a damped oscillatory gravitational-wave signature. TFT reinterprets the mechanical origin of that signal as relaxation of a distorted saturated boundary, rather than treating it only as a geometric perturbation of the external black-hole spacetime.

In the full nonlinear merger problem, different deformation modes of the boundary would generally have different damping rates and relaxation frequencies. The single equation above is therefore not a complete waveform model. It is the lowest-order mechanical closure showing how merger distortion, stress shedding, and ringdown are connected in the live phase-boundary framework.

\section{Kinematic Capacity Constraint}

The Soft Floor has a finite elastic kinematic capacity. A live black-hole boundary therefore cannot support unlimited simultaneous motion in all of its dynamical channels. Center-of-mass translation, phase-front growth, tangential redistribution, and torsional boundary motion all draw on the same underlying capacity of the surrounding FCC vacuum to deform, transmit stress, and remain active.

In a first-order channel approximation, this may be written schematically as:
\begin{equation}
    V_{\mathrm{CM}}^2 + v_n^2 + v_{\parallel}^2 + v_{\mathrm{tor}}^2 \lesssim c_{\mathrm{E}}^2.
\end{equation}

Here $V_{\mathrm{CM}}$ is the center-of-mass speed of the saturated black-hole domain, $v_n$ is the normal phase-front speed, $v_{\parallel}$ is the tangential boundary-redistribution speed, and $v_{\mathrm{tor}}$ is the effective torsional boundary speed associated with axial angular momentum.

This expression should not be read as a completed microscopic field equation. It is a mechanical capacity rule. Its purpose is to state that the live FCC/SC boundary must distribute the available Soft Floor response among different modes of motion.

When the combined dynamical demand remains below the Soft Floor capacity, the boundary can respond through smooth orbital motion, local growth, redistribution, and torsional loading. When the demand approaches the capacity scale $c_{\mathrm{E}}^2$, smooth boundary response is no longer sufficient. The system must instead respond through one or more nonlinear mechanisms:
\begin{equation}
    \text{phase locking}, \qquad \text{common-boundary formation}, \qquad \text{stress shedding}, \qquad \text{surface relaxation}.
\end{equation}

This connects binary inspiral, saturated bridge formation, gravitational-wave emission, and ringdown into a single mechanical sequence. In TFT, merger is not merely an orbital event; it is the point at which the live boundary and intervening Soft Floor can no longer maintain separate active FCC response under the combined strain load.

A more complete mechanism would replace the schematic scalar condition with a tensorial surface-capacity law, including coupling between normal growth, tangential shear, torsion, and external Soft Floor strain. For the present appendix, the capacity condition serves only as the minimal mechanical closure needed to keep the velocity channels finite and mutually constrained.

\section{Observational Discipline}

This appendix does not claim that TFT has already replaced the Kerr solution, numerical-relativity merger simulations, or gravitational-wave waveform models. The claim is narrower.

TFT provides a mechanical interpretation and first-order surface-dynamics framework for black-hole spin, binary merger, stress shedding, and ringdown. It identifies the live FCC/SC phase boundary as the mechanical carrier of the dynamics that standard black-hole physics represents geometrically.

The present appendix therefore establishes the following interpretive map:
\begin{itemize}
    \item $\text{black-hole spin} \rightarrow \text{torsional boundary motion and external Soft Floor shear}$,
    \item $\text{binary merger} \rightarrow \text{saturated bridge formation and common FCC/SC boundary}$,
    \item $\text{gravitational-wave emission} \rightarrow \text{stress shedding from rapid boundary reconfiguration}$,
    \item $\text{ringdown} \rightarrow \text{damped relaxation of the distorted live phase boundary}$.
\end{itemize}

A full TFT replacement of strong-field black-hole dynamics would still require future derivation of:
\begin{equation}
    \text{strong-field boundary motion, surface torsional stress, frame-dragging analogue},
\end{equation}
\begin{equation}
    \text{ISCO analogue, waveform generation, ringdown frequency spectrum}.
\end{equation}

The purpose of this appendix is therefore not to present a completed waveform theory. Its purpose is to define the mechanical degrees of freedom that such a theory would require.

The contribution may be summarized as:
\begin{equation}
    \boxed{\text{TFT converts black-hole dynamics from pure metric geometry into physical live phase-boundary mechanics.}}
\end{equation}

\section{Conclusion}

This appendix extends TFT black-hole physics from spherical saturation to live surface dynamics.

A TFT black hole has center-of-mass motion, orbital motion, axial angular momentum, phase-front growth, merger dynamics, gravitational-wave stress shedding, and ringdown. These dynamics are not assigned to ordinary material flow through the saturated SC interior. They are assigned to the FCC/SC phase boundary and to the surrounding strained Soft Floor.

The saturated SC domain remains internally phase-locked. Observed black-hole spin is therefore not modeled as ordinary fluid-like circulation inside the core. In TFT, axial angular momentum is carried by torsional loading of the live boundary and by shear strain in the external Soft Floor.

Binary black-hole merger is not a zero-distance collision. It is saturated bridge formation. Two initially separate saturated domains become one larger phase-boundary system when the intervening FCC Soft Floor reaches the saturation condition and locks into the SC phase:
\begin{equation}
    \alpha_{\mathrm{bridge}} = \alpha_{\mathrm{G},1} + \alpha_{\mathrm{G},2} + \alpha_{\mathrm{dyn}} \geq 1.
\end{equation}

The ideal circular boundary-contact limit gives $v_{\mathrm{rel}} = c_{\mathrm{E}}$, but only at the contact scale $d = R_1 + R_2$. The general separation-controlled relation is:
\begin{equation}
    v_{\mathrm{rel}} = c_{\mathrm{E}} \sqrt{\frac{R_1+R_2}{d}}.
\end{equation}
Thus, $c_{\mathrm{E}}$ is a contact-limit velocity scale, not a universal orbital speed for all binary separations.

Gravitational waves are interpreted as stress-shedding waves produced by rapid reconfiguration of the live FCC/SC boundary. Ringdown is the damped relaxation of the distorted post-merger phase boundary:
\begin{equation}
    \text{ringdown} = \text{damped relaxation of the live FCC/SC phase boundary}.
\end{equation}

The minimum mechanical structure required for TFT black-hole dynamics is therefore:
\begin{equation}
\begin{split}
    &\text{mass-bearing boundary} + \text{phase-front growth} + \text{torsional spin} \\
    &+ \text{saturated bridge merger} + \text{stress-shedding radiation}.
\end{split}
\end{equation}

This appendix establishes that structure as a first-order surface theory. It does not complete the full rotating or binary black-hole solution, but it defines the mechanical program required for TFT to develop one.